\section{Mathematical model}\label{sect:Model}


Consider a two-dimensional liquid droplet, as though an infinite-cylinder in the out-of-plane direction, moving through an air-filled upper half plane towards a deep bath of infinite depth and width as shown in figure \ref{fig:figure1}. We restrict the centre of mass of the droplet to vertical motion along the $z$-axis, introducing a symmetry to the problem with respect to the $z$ axis. Introducing the subscript notation $\cdot_\alpha$ to denote the different fluid regions, where $\alpha \in \{a,b,d\}$ are the air, bath, and droplet respectively, the velocity of each fluid is given by $\mathbf{u}_\alpha$, the pressure is given by $p_\alpha$, and $\mu_\alpha$, $\rho_\alpha$ and  $g$ denote the dynamic viscosity, density and gravity, respectively.

\begin{figure}
    \centering
    \include{Figures/Fig1}
    \caption{Schematic of model domain denoting two frame of reference points. The origin for the full frame is the Cartesian $(x,z)$-coordinate set directly beneath the south pole of the drop and aligned to an unperturbed bath free surface. The droplet will be evolved in a moving frame of reference with polar coordinates ($r,\theta)$ about the centre of mass of a drop whose free surface lies at $r=R_0$, and $\theta=0$ pointing vertically downwards towards the bath. }
    \label{fig:figure1}
\end{figure}
While each region of fluid will be treated differently due to particular approximations and geometries, in each region, $\mathbf{\Omega}_\alpha$, we begin with the incompressible Navier-Stokes equations

   \begin{align}\label{eq: mass}
    \nabla \cdot \mathbf{u}_\alpha &=  0, \qquad &\mathbf{x} &\in \Omega_\alpha,\\ \label{eq: mom}
    \rho_\alpha \left(\partial_t \mathbf{u}_\alpha + (\mathbf{u}_\alpha\cdot \nabla)\mathbf{u}_\alpha \right) &= -\nabla p_\alpha + \mu_\alpha \nabla^2 \mathbf{u}_\alpha - \rho_\alpha g \mathbf{e}_z, \qquad & \mathbf{x} &\in \Omega_\alpha,
\end{align}
where $g$ is the gravitational constant.
At the interface between fluids, defined by $F_\beta(\mathbf {x},t)=0$, where $\beta$ takes on the values $b$ or $d$, for the bath-air and drop-air interfaces respectively, the kinematic and stress continuity boundary conditions and the velocity continuity respectively are given by 

\begin{align}\label{eq: KBC General}
    \partial_t F_\beta + (\mathbf{u}_\beta \cdot \nabla)F_\beta &= 0, \qquad &x\in \partial\Omega_\beta,\\
    p_\beta \mathbf{n}_\beta - \mu_\beta \mathbf{\tau}_\beta \mathbf{n}_\beta &= p_a \mathbf{n}_\beta - \mu_a \mathbf{\tau}_a  \mathbf{n}_\beta + \sigma_\beta \kappa_\beta \mathbf{n}_\beta, \qquad & x\in \partial\Omega_\beta,\label{eq:StressCont} \\
    \mathbf{u}_\beta &=  \mathbf{u}_a,  \qquad & x\in \partial\Omega_\beta,\label{eq:VelocityCont}
   \end{align}
   
where $\sigma$ is the surface tension coefficient, $\kappa$ denotes the curvature of the free surface whose unit normal vector is given as $\mathbf{n}_\beta$, and  $\tau = (\nabla \mathbf{u} + \nabla \mathbf{u}^T)$ is the rate of strain tensor. In the far field for the air and the bath, the velocities vanish, and pressures tend to their hydrostatic balance values. 
We shall now consider the flows within the bath, air, and liquid domains separately, using the relevant approximation the respective domain requires.
   
\subsection{Linear quasi-potential bath model}


Following methodology well established in the literature \cite{milewski2015faraday,galeano2017non,galeano2021capillary, alventosa2023inertio} to describe the waves generated by a bath-droplet interaction, we proceed by assuming a fluid bath with small free-surface perturbations from the equilibrium $z=0$  denoted by $z=\eta_b (x,t)$. As the liquid is incompressible we can take the velocity field of the bath as  $\mathbf{u}_b = (u_b, w_b) = \nabla \phi + \nabla^\perp \psi$, where $\nabla \phi$ denotes the irrotational velocity component, and $\psi$ a small vortical correction arising from viscous effects. Here we have used the notation $\nabla^\perp = (-\partial_z, \partial_x)$. Due to the assumption that disturbances are small, the system \eqref{eq: mass} -- \eqref{eq:VelocityCont} results in the following linearized problem

\begin{align}
\nabla^2 \phi &= 0, \qquad &z&\le 0, \label{eq: laplace}\\
\rho_b \partial_t \phi &= - p_b - \rho_b g z, \qquad &z&\le 0, \\
\rho_b \partial_t \psi &= \mu_b \nabla^2 \psi, \qquad &z& \le 0, 
\end{align}
and boundary conditions

\begin{align}
\partial_t \eta_b  &= \partial_z \phi + \partial_x \psi , \qquad &z&=0,  \\
 p_b - p_a - \sigma_b \kappa_b &= 2\mu_b\left(\partial_z^2 \phi + \partial_{xz} \psi \right), \qquad &z&=0, \\
   \mu_a (\partial_z u_a + \partial_x w_a) &= \mu_b (2\partial_{xz} \phi + \partial_x^2 \psi - \partial_z^2 \psi) , \qquad &z&=0, \\
   u_a &= \partial_{x} \phi - \partial_z \psi, \qquad &z&=0, \\
    w_a &= \partial_{z} \phi + \partial_x \psi, \qquad &z&=0, \label{eq: z velocity match} \\
\phi,\psi &\to 0, \qquad &z &\to -\infty, \label{eq: end of full system}
\end{align}
which include the kinematic boundary condition, normal and tangential stress balances, velocity continuity, and far field decay, respectively. We have omitted normal stresses in the air, in anticipation of a lubrication layer approximation.

The arguments in Dias \textit{et al} \cite{dias2008theory} (building on the viscous damping rates for water waves obtained by Lamb \cite{lamb1924hydrodynamics}) involve a smallness assumption in the boundary layer flow represented by $\psi$ and a boundary layer scaling for $\psi$ in its vertical dependence. This is explained in more detail in Appendix \ref{App: Boundary}. The boundary conditions then simplify to

\begin{align}
\partial_t \eta_b  &= \partial_z \phi + \partial_x \psi , \qquad &z&=0,  \\
\partial_t \phi &=  - \frac{1}{\rho_b} p_a + g \eta_b + \frac{\sigma_b}{\rho_b} \partial_x^2 \eta_{b} - 2\nu_b \partial_z^2 \phi, \qquad &z&=0, \\
    \partial_t \psi &=  \nu_b \left(2+\frac{\mu_a}{ \mu_b}\right) \partial_{zx} \phi + \nu_b \frac{\mu_a}{\mu_b}\partial_z u_a , \qquad &z&=0, \\
   u_a &= \partial_{x} \phi, \qquad &z&=0, \\
    w_a &= \partial_{z} \phi, \qquad &z&=0,
\end{align}
where we have introduced  the kinematic viscosity $\nu_b = \mu_b/\rho_b$ and substituted for the linearized curvature of the surface $\kappa = -\partial_x^2 \eta_b$. Disregarding terms of $O({\mu_a}/{\mu_b})$, one can deduce using the leading order balance $\partial_t \eta_b = \partial_z \phi$ that $\partial_x \psi = 2\nu_b \partial_{x}^2 \eta_b$, and hence the boundary conditions for the velocity potential and the free surface evolution are forced by air pressure alone

\begin{align}
\partial_t \phi &= -\frac{1}{\rho_b} p_a -  g\eta_b + \frac{\sigma_b}{\rho_b}\partial_x^2 \eta_b  + 2\nu_b \partial_x^2 \phi, \qquad &z&=0, \label{eq: qp DBC} \\
\partial_t \eta_b  &= \partial_z \phi + 2\nu_b \partial_x^2 \eta_b, \qquad &z&=0, \label{eq: qp KBC}
\end{align}
where $\phi$ satisfies Laplace's equation in the lower half plane. The conditions for velocity continuity at the interface will be matched in the equations for the lubrication air layer, and the pressure $p_a$ will be taken to be zero outside this lubrication layer. We note that the kinematic boundary condition requires $\partial_z \phi$ to be evaluated at $z=0$, and that this can be expressed simply in terms of $\phi$ at the surface once the system is Fourier transformed in $x$ as shown in Section \ref{sect:Numerics}.
\subsection{Lubricating air layer model}

When the droplet of undisturbed radius $R_0$ is sufficiently far from the bath, the two liquid regions have a negligible effect on each other. However, when the height between the drop and the bath is of $\mathcal{O}(\varepsilon R_0)$, where $\varepsilon \ll 1$, they can start to affect each other and we are able to proceed by considering the air layer and the pressure within as a lubrication region \cite{hicks2013liquid}. 

The geometry of this lubrication region is given by
\begin{equation*}
\Omega_{l^*}(t) = \left\{(x,z): x \in [-l^*(t),l^*(t)], z \in [\eta_b(x,t),\eta_d(x,t)] \right\},    
\end{equation*}
where $l^*$ is the horizontal location at which the vertical distance between the bath and droplet becomes greater than $\varepsilon R_0$ for some chosen $\varepsilon \ll 1$. In the next section when we consider the evolution of the droplet it will be more appropriate to set $\eta_d = Z(t)-S(x,t)$, where $Z(t)$ is the vertical height of the centre of mass of the drop, and $S(x,t)$ is the shape profile of the lower interface of the droplet relative to that centre of mass. For rigid objects $S$ is assumed to be constant in time. The height of the lubrication layer within its horizontal domain, is given by 
\begin{equation}
    h(x,t) = \eta_d - \eta_b.
\end{equation}

Taking this geometry into consideration, and assuming weak impacts in which the lubrication layer will not deviate substantially from the horizontal, \eqref{eq: mass}-\eqref{eq: mom} can be reduced  with lubrication scaling to

\begin{align}
    \partial_x u_a + \partial_z w_a &= 0\label{eq: air mass}, \\
     0 &= -\partial_x p_a + \mu_a \partial_z^2 u_a,\label{eq: uzz} \\
    0 &= -\partial_z p_a. \label{eq: pz}
\end{align}
From equation \eqref{eq: pz}, we see that the pressure within the air layer varies in the horizontal direction only, i.e $p_a = P(x,t)$. At the two air-liquid interfaces we impose velocity matching as in \eqref{eq:VelocityCont}, such that

\begin{equation*}
    z=\eta_b: \quad u_a = \partial_x \phi, \quad w_a  =\partial_z \phi;  \qquad  \qquad z=\eta_d: \quad
    u_a = u_d, \quad w_a = w_d,
\end{equation*}
where we recall that $(u_d,w_d)$ are the velocity components of the flow in the droplet. Equation \eqref{eq: uzz} can then be solved to determine the form of $u_a$ as a Poiseuille flow within the air layer of the form 

\begin{equation}\label{eq: u_a}
    u_a = \frac{\partial_x P}{2\mu_a}(z-\eta_b)(z-\eta_d) + \frac{u_d|_{\eta_d}}{\eta_d- \eta_b}(z-\eta_b) + \frac{\partial_x \phi_b|_{\eta_b}}{\eta_b- \eta_d}(z-\eta_d).
\end{equation}

The lubrication equation results from integrating the conservation of mass equation \eqref{eq: air mass} across the air layer, 

\begin{equation}
    \int_{\eta_b}^{\eta_d} \partial_x u_a + \partial_z w_a \; \text{d}z = 0.
\end{equation}
Using the Leibniz integral rule 

\begin{equation}
    \partial_x \int_{\eta_b}^{\eta_d} u_a \text{d}z + \left(w_a\Big|_{\eta_d} - u_a \Big|_{\eta_d}\partial_x \eta_d\right)  - \left(w_a\Big|_{\eta_b} - u_a \Big|_{\eta_b} \partial_x \eta_b \right) = 0,
\end{equation}
which is a conservation law for the height of the layer, since the first term represents the area flux, and the boundary terms yield, from \eqref{eq: KBC General}, the change in thickness i.e. $\partial_x Q + \partial_t \eta_d  - \partial_t \eta_b = 0$ with a flux from (\ref{eq: u_a})

\begin{equation}
Q = -\frac{\partial_x P}{12\mu_a} h^3 + \frac{(\partial_x \phi_b|_{\eta_b} + u_d|_{\eta_d})}{2} h~.
\end{equation}
Hence the lubrication equation is

\begin{equation}\label{eq: thin film}
\partial_t h +  \partial_x \left[-\frac{1}{12\mu_a} \partial_x P h^3 + \frac{(\partial_x \phi_b|_{\eta_b} + u_d|_{\eta_d} )}{2} h \right] = 0.
\end{equation}
The boundary conditions for this equation are 
$P=0$ at $x=\pm l^*$.
\subsection{Droplet behaviour}

\subsubsection{Droplet trajectory}
The vertical motion of the droplet can be described using its centre of mass $Z(t)$. We shall ignore the air effects (primarily viscous drag) except for the interaction with the lubrication layer. Taking the mass of the droplet to be given by $\rho_d \pi R_0^2$, where $R_0$ is the undeformed radius of the drop, we write the vertical acceleration of the droplet as 

\begin{equation}\label{eq: vert dyn}
    \rho_d \pi R^2 \frac{\text{d}^2 Z}{\text{d}t^2} = \int_{-l^*}^{l^*} P dx - \rho_d \pi R^2 g,
\end{equation}
where the integral is the force due to the pressure from the thin film region, and $g$ is the gravitational constant.

If the droplet were to be modelled by a non-deformable disk we can take $h = \eta_d - \eta_b = (Z(t)-S(x))-\eta_b(x,t)$, with $u_d|_{\eta_d}=0$, and the system of equations consisting of (\ref{eq: laplace}), (\ref{eq: qp KBC}), (\ref{eq: qp DBC}), (\ref{eq: thin film}), and (\ref{eq: vert dyn}) is now closed.

\subsubsection{Droplet deformation}
The equations which govern the evolution of drop deformations follow the methodology first used by Lamb \cite{lamb1924hydrodynamics}, and more recently described by Aalilija \textit{et al} \cite{aalilija2020analytical} and Alventosa \textit{et al} \cite{alventosa2023inertio}. 
We describe the motion of fluid within the droplet in polar coordinates, with origin at the centre of mass of the drop and $\theta =0$ in the negative $z$-direction. The free surface of the droplet at rest lies along $r=R_0$, with small perturbations to this given by $\eta_d(\theta,t)$. 

In order to derive the equations one can re-cast the equations (\ref{eq: laplace})-(\ref{eq: z velocity match}) in polar coordinates, and proceed with a reduction for the viscous boundary layer near $r=R_0$ to obtain quasi-potential boundary conditions corresponding to (\ref{eq: qp KBC}) - (\ref{eq: qp DBC}). The resulting problem to solve is 

\begin{align}
 \nabla^2 \phi_d &= 0, \quad &r&<R_0 \\
    \partial_t \phi_d &= -\frac{1}{\rho_d}p_a - \frac{\sigma}{\rho_d} \kappa  - 2\nu_d\partial_r^2 \phi_d, \quad &r&=R_0 \label{eq: drop DBC}\\
    \partial_t \eta_d &= \partial_r \phi_d + 2 {\nu_d}\left(\frac{1}{r^2} \partial_\theta^2\eta_d - \frac{1}{r^3}\partial_\theta^2 \int \phi_d \;\text{d} t\right), \quad &r&=R_0, \label{eq: drop KBC}
\end{align}
where $\kappa = -\frac{1}{R^2_0}\left(\eta_d + \partial_\theta^2 \eta_d\right)$ is the linearized perturbation of curvature in polar coordinates. Solving Laplace's equation and expanding free surface deformations and pressure in terms of Fourier series symmetric about $\theta=0$, we require that

\begin{align}
    \phi_d(r, \theta, t) &= \sum_{n=2}^\infty a_n(t)r^n cos(n\theta), \\
    \eta_d(\theta, t) &= \sum_{n=2}^\infty c_n(t) cos(n\theta), \\
    p_a &= \sum_{n=2}^\infty \hat{p}_n(t) cos(n\theta),
\end{align}

where the $n=0$ terms are omitted to enforce mass conservation, and the $n=1$ terms are omitted as they correspond to the linearized center of mass evolution. Substitution into the expressions (\ref{eq: drop DBC})-(\ref{eq: drop KBC}) results in

\begin{align}
    \dot{a}_n R_0^n &= -\frac{1}{\rho_d}\hat{p}_n + \frac{\sigma}{\rho R_0^{2}} \left( 1-n^2\right)c_n - 2\nu_d \frac{1}{R_0^{2}}n(n-1)a_n R_0^n,\\
    \dot{c}_n &= \frac{1}{R_0} n a_n R_0^{n} + 2\nu_d \left(-\frac{1}{R_0^2}n^2 c_n  + \frac{1}{R_0^{3}} n^2  \int a_n R_0^n  \;\text{d}t\right). 
\end{align}
 
Considering the leading order balance $\dot{c}_n = \frac{1}{R_0} n a_n R_0^{n}$ into the viscous terms, and eliminating coefficients $a_n$ by differentiating the second equation and using the first equation, results in the evolution equation for the Fourier coefficient of the free surface shape 

\begin{equation}
    \ddot{c}_n + 2\lambda_n \dot{c}_n + \omega_n^2 c_n = -\frac{n}{\rho_d R_0}\hat{p}_n,
\end{equation}
where 

\begin{equation}
    \lambda_n = \frac{2 \mu_d}{\rho_d R_0^2}n(n-1), \qquad \omega_n^2 = \frac{\sigma_d }{\rho_d R_0^3}n(n^2-1) 
\end{equation}
are viscous damping and frequency terms, respectively. We note that these correspond to the values in Aalilija \textit{et al}\cite{aalilija2020analytical} using the energy method of Lamb\cite{lamb1924hydrodynamics}, but were obtained directly from the quasi-potential approximation. The surface tangential velocity which appears in the lubrication equation (\ref{eq: thin film}) can be recovered by the leading order 

\begin{equation}
    u_d(R_0, \theta, t) = -\sum_{n=2}^\infty \dot{c}_n(t) sin(n\theta).
\end{equation}

We have now closed the system of equations which describe fully the lubrication-mediated behaviour of a droplet rebounding off a deep liquid bath. Each of the three regions of fluid have been prescribed a governing set of equations, which are all coupled together through the presence of the pressure function $P(x,t)$. We will now present a numerical implementation of the system in Section \ref{sect:Numerics}.



